\section{Map small signed differences to small unsigned integers}
A two's-complement negative integer has high one bits. Treating that full bit pattern
as an unsigned variable-length value makes even \(-1\) large. Use the following
signed-to-unsigned mapping, often called ZigZag:
\[
  u=\begin{cases}2d&d\geq0,\\-2d-1&d<0.\end{cases}
\]
Nonnegative deltas produce even values and negative deltas produce odd values.
The low bit records the sign choice while the remaining bits record magnitude.
The inverse is
\[
  d=\begin{cases}u/2&u\text{ even},\\-(\lfloor u/2\rfloor)-1&u\text{ odd}.\end{cases}
\]

\begin{table}[htbp]
  \centering
  \caption{Nearby signed deltas and their unsigned mapped values.}
  \tagpdfsetup{table/header-rows={1}}
  \begin{tabular}{rrrrrrr}
    \toprule
    Delta  & \(-3\) & \(-2\) & \(-1\) & 0 & 1 & 2 \\
    \midrule
    Mapped & 5      & 3      & 1      & 0 & 2 & 4 \\
    \bottomrule
  \end{tabular}
\end{table}

For a 64-bit register containing the exact signed delta, the same mapping is
\((d\ll1)\mathbin{\mathrm{xor}}(d\gg63)\), with an arithmetic right shift.
The right shift yields either zero or a full sign mask. Use \texttt{dsll} and
\texttt{dsra}; the latter must shift by 63, not by a reading's stored width minus
one. A positive exact delta can exceed signed-word range. Its upper bits are not a
sign mask obtained by shifting only thirty-one places.

The largest mapped value an encoder can produce is \(8589934590=2^{33}-2\).
Keep all 33 relevant bits. The decoder's wire parser accepts the bounded unsigned
range through \(2^{33}-1\), then checks whether its inverse delta and predecessor
produce a legal reading. Wire width and reading range are separate validations.

\begin{MARSExercise}{Invert a mapped value}
  Decode mapped values 0, 1, 2 and 3 without using a table. Which parity selects a
  negative delta? Then compute the inverse of \(2^{33}-1\) and determine whether any
  permitted predecessor can make that delta reconstruct a signed 32-bit reading.
\end{MARSExercise}
